\chapter{React Project Structure}

A flat \texttt{src/} folder works for a demo, but breaks down once an app has
auth, protected routes, and multiple pages. Here is the structure used
throughout this book for the Task Management System.

\begin{diagrambox}[frontend/src Structure]
\begin{verbatim}
frontend/src/
|-- api/axios.js          (single configured axios instance)
|-- services/             authService.js, taskService.js
|-- components/           TaskCard.jsx, TaskList.jsx, Navbar.jsx
|-- pages/                Dashboard.jsx, Login.jsx, TaskDetails.jsx
|-- layouts/MainLayout.jsx (Navbar + Sidebar + <Outlet/>)
|-- hooks/useAuth.js
|-- context/AuthContext.jsx
|-- utils/formatDate.js
|-- assets/logo.svg
|-- App.jsx
`-- main.jsx
\end{verbatim}
\end{diagrambox}

\section{What Goes Where}

\begin{table}[h]
\centering
\begin{tabularx}{\textwidth}{lXX}
\toprule
\textbf{Folder} & \textbf{Contains} & \textbf{Should NOT contain} \\
\midrule
\texttt{components/} & Reusable UI pieces & Routes, data fetching \\
\texttt{pages/} & Route-level screens & Shared low-level UI \\
\texttt{layouts/} & Shell wrapping pages & Feature-specific logic \\
\texttt{hooks/} & Reusable stateful logic & JSX markup \\
\texttt{context/} & Global providers/state & One-off local state \\
\texttt{services/} & Resource-specific API calls & Raw axios config, JSX \\
\texttt{api/} & Axios instance, interceptors & Resource business logic \\
\texttt{utils/} & Pure helper functions & Components, hooks \\
\bottomrule
\end{tabularx}
\end{table}

\begin{notebox}[NOTE]
\texttt{api/axios.js} is generic HTTP plumbing; \texttt{services/taskService.js}
knows the specific endpoints. This split means a backend URL change never
touches business logic.
\end{notebox}

\begin{tipbox}[TIP]
Quick test: needs its own route in \texttt{App.jsx}? It's a page. Used in two
or more pages? It's a component.
\end{tipbox}
